\documentclass[10pt,a4paper]{amsart}
\usepackage{amsmath,amssymb,amsthm,graphicx}
% Retained minimal preamble for the Stacks Project chapter flat.tex (More on Flatness).
% Verbatim excerpt of the chapter preamble: Xy-pic setup (native preamble lines 14-18) and the native
% theorem declarations with the mathematical macros (native preamble lines 154-217). Copied mechanically;
% nothing added, removed or redefined. The wrapper class and the amsmath/amssymb/amsthm/graphicx base
% packages are supplied by the assembler, so the class/externaldocument part of the native preamble is omitted.
% For commutative diagrams we use Xy-pic
\usepackage[all]{xy}

% We use 2cell for 2-commutative diagrams.
\xyoption{2cell}
% Theorem environments.
%
\theoremstyle{plain}
\newtheorem{theorem}[subsection]{Theorem}
\newtheorem{proposition}[subsection]{Proposition}
\newtheorem{lemma}[subsection]{Lemma}

\theoremstyle{definition}
\newtheorem{definition}[subsection]{Definition}
\newtheorem{example}[subsection]{Example}
\newtheorem{exercise}[subsection]{Exercise}
\newtheorem{situation}[subsection]{Situation}

\theoremstyle{remark}
\newtheorem{remark}[subsection]{Remark}
\newtheorem{remarks}[subsection]{Remarks}

\numberwithin{equation}{subsection}

% Macros
%
\def\lim{\mathop{\mathrm{lim}}\nolimits}
\def\colim{\mathop{\mathrm{colim}}\nolimits}
\def\Spec{\mathop{\mathrm{Spec}}}
\def\Hom{\mathop{\mathrm{Hom}}\nolimits}
\def\Ext{\mathop{\mathrm{Ext}}\nolimits}
\def\SheafHom{\mathop{\mathcal{H}\!\mathit{om}}\nolimits}
\def\SheafExt{\mathop{\mathcal{E}\!\mathit{xt}}\nolimits}
\def\Sch{\mathit{Sch}}
\def\Mor{\mathop{\mathrm{Mor}}\nolimits}
\def\Ob{\mathop{\mathrm{Ob}}\nolimits}
\def\Sh{\mathop{\mathit{Sh}}\nolimits}
\def\NL{\mathop{N\!L}\nolimits}
\def\CH{\mathop{\mathrm{CH}}\nolimits}
\def\proetale{{pro\text{-}\acute{e}tale}}
\def\etale{{\acute{e}tale}}
\def\QCoh{\mathit{QCoh}}
\def\Ker{\mathop{\mathrm{Ker}}}
\def\Im{\mathop{\mathrm{Im}}}
\def\Coker{\mathop{\mathrm{Coker}}}
\def\Coim{\mathop{\mathrm{Coim}}}

% Boxtimes
%
\DeclareMathSymbol{\boxtimes}{\mathbin}{AMSa}{"02}

%
% Macros for moduli stacks/spaces
%
\def\QCohstack{\mathcal{QC}\!\mathit{oh}}
\def\Cohstack{\mathcal{C}\!\mathit{oh}}
\def\Spacesstack{\mathcal{S}\!\mathit{paces}}
\def\Quotfunctor{\mathrm{Quot}}
\def\Hilbfunctor{\mathrm{Hilb}}
\def\Curvesstack{\mathcal{C}\!\mathit{urves}}
\def\Polarizedstack{\mathcal{P}\!\mathit{olarized}}
\def\Complexesstack{\mathcal{C}\!\mathit{omplexes}}
% \Pic is the operator that assigns to X its picard group, usage \Pic(X)
% \Picardstack_{X/B} denotes the Picard stack of X over B
% \Picardfunctor_{X/B} denotes the Picard functor of X over B
\def\Pic{\mathop{\mathrm{Pic}}\nolimits}
\def\Picardstack{\mathcal{P}\!\mathit{ic}}
\def\Picardfunctor{\mathrm{Pic}}
\def\Deformationcategory{\mathcal{D}\!\mathit{ef}}

% --- appended 2026-10-12 (Chief): the project own definitions of the reference/slogan/history
% display environments, copied mechanically from the shared Stacks preamble lines 8-12
% (\usepackage{verbatim} plus the three \newenvironment lines, which the chapter preamble keeps
% verbatim from the project preamble). Nothing added, removed or redefined: these environments
% carry no mathematical content and the project itself typesets them as comments. Without these
% three lines a counted statement whose own text uses them fails with Environment slogan undefined.
\usepackage{verbatim}
\newenvironment{reference}{\comment}{\endcomment}
\newenvironment{slogan}{\comment}{\endcomment}
\newenvironment{history}{\comment}{\endcomment}
\usepackage{cleveref}
\crefname{originaldefinition}{original definition}{original definitions}
\Crefname{originaldefinition}{Original definition}{Original definitions}
\makeatletter
\newlabel{definition-standard-h}{{0EUY}{source}{}{sourcepointer-1}{}}
\newlabel{definition-standard-h@cref}{{[originaldefinition][0][]0EUY}{[1][0][]source}{}{}{}}
\makeatother
\crefname{originallemma}{original lemma}{original lemmas}
\Crefname{originallemma}{Original lemma}{Original lemmas}
\makeatletter
\newlabel{limits-lemma-descend-finite-presentation}{{01ZM}{source}{}{sourcepointer-2}{}}
\newlabel{limits-lemma-descend-finite-presentation@cref}{{[originallemma][0][]01ZM}{[1][0][]source}{}{}{}}
\makeatother
\crefname{originallemma}{original lemma}{original lemmas}
\Crefname{originallemma}{Original lemma}{Original lemmas}
\makeatletter
\newlabel{limits-lemma-descend-affine-finite-presentation}{{01ZN}{source}{}{sourcepointer-3}{}}
\newlabel{limits-lemma-descend-affine-finite-presentation@cref}{{[originallemma][0][]01ZN}{[1][0][]source}{}{}{}}
\makeatother
\crefname{originallemma}{original lemma}{original lemmas}
\Crefname{originallemma}{Original lemma}{Original lemmas}
\makeatletter
\newlabel{lemma-approximate-h-cover}{{0ETU}{source}{}{sourcepointer-4}{}}
\newlabel{lemma-approximate-h-cover@cref}{{[originallemma][0][]0ETU}{[1][0][]source}{}{}{}}
\makeatother
\crefname{originallemma}{original lemma}{original lemmas}
\Crefname{originallemma}{Original lemma}{Original lemmas}
\makeatletter
\newlabel{limits-lemma-eventually-proper}{{081F}{source}{}{sourcepointer-5}{}}
\newlabel{limits-lemma-eventually-proper@cref}{{[originallemma][0][]081F}{[1][0][]source}{}{}{}}
\makeatother
\crefname{originallemma}{original lemma}{original lemmas}
\Crefname{originallemma}{Original lemma}{Original lemmas}
\makeatletter
\newlabel{limits-lemma-descend-surjective}{{07RR}{source}{}{sourcepointer-6}{}}
\newlabel{limits-lemma-descend-surjective@cref}{{[originallemma][0][]07RR}{[1][0][]source}{}{}{}}
\makeatother
\crefname{originallemma}{original lemma}{original lemmas}
\Crefname{originallemma}{Original lemma}{Original lemmas}
\makeatletter
\newlabel{limits-lemma-descend-opens}{{01Z4}{source}{}{sourcepointer-7}{}}
\newlabel{limits-lemma-descend-opens@cref}{{[originallemma][0][]01Z4}{[1][0][]source}{}{}{}}
\makeatother
\title{The category of standard h coverings of an affine limit is the colimit of the categories of standard h coverings of the approximating schemes}
\author{The Stacks Project Authors}
\date{stacks-project:flat.tex:local-snapshot-20261011; GNU FDL 1.2 or later}
\begin{document}
\maketitle
\noindent\textit{Editorial note.} All delivered mathematical text is the author's own, copied from the pinned Stacks Project chapter file flat.tex; no statement or proof step was rewritten, reordered or removed.\par\medskip
\noindent\textit{Editorial note.} The preamble is the chapter's own retained minimal preamble (sources/stacks-flat/retained-preamble.tex), a verbatim excerpt of the native Xy-pic setup and of the native theorem declarations and mathematical macros, followed by the project own definitions of the reference, slogan and history display environments (verbatim package plus three newenvironment lines) that the retained file carries at its end and that the project itself typesets as comments. No macro was invented or redefined, and the wrapper class and the amsmath/amssymb/amsthm/graphicx base packages are supplied by the assembler.\par\medskip
\noindent\textit{Editorial note.} The counted statement and proof contain no cross reference, no citation and no included file, so the delivered body is self contained.\par\medskip
\noindent\textit{Editorial note.} Canonical location of the counted proposition: Stacks tag 0EV2 (sources/stacks-flat/pinned-original-source.tex line 10768).\par\medskip
\noindent\textit{Allowed standard prerequisites (editorial).}\begin{itemize}
\item Standard h coverings and the definition of an h covering. Same chapter, Definition labelled definition-standard-h, cited in the statement. Fixes the objects whose category is being approximated.
\item Approximating an h covering by a proper finitely presented morphism with an affine cover. Same chapter, Lemma labelled lemma-approximate-h-cover, cited at proof lines 10809-10811. Replaces the h covering by a finitely presented proper surjective morphism with a refining affine cover.
\item Descent of finite presentation, affineness and of morphisms of finite presentation through a directed limit. Limits, Lemmas labelled limits-lemma-descend-finite-presentation and limits-lemma-descend-affine-finite-presentation, cited at proof lines 10797-10801. Identifies the categories of finitely presented (affine) schemes over T with the colimit of those over the T\_i.
\item Descending properness, surjectivity and opens. Limits, Lemmas labelled limits-lemma-eventually-proper, limits-lemma-descend-surjective and limits-lemma-descend-opens, cited at proof lines 10815-10821. Descends the approximating proper cover and its affine open cover to some T\_\{i'\}.
\end{itemize}\medskip
\ifdefined\hypertarget\hypertarget{sourcepointer-1}{}\fi
\noindent\textit{Source reference.} Original definition 0EUY: sources/stacks-flat/pinned-original-source.tex, source lines 10462--10462. Same-chapter pointer into the author environment definition at source line 10462; invoked at proof line None of the counted proof. It is retained only as a navigation pointer; every load-bearing step is inside the delivered span.\par\medskip
\ifdefined\hypertarget\hypertarget{sourcepointer-2}{}\fi
\noindent\textit{Source reference.} Original lemma 01ZM: sources/stacks-chow/tags.txt, source lines 2599--2599. Cross-chapter standard prerequisite located in the project tag list staged in the local snapshot; invoked at proof line 10798 of the counted proof. It is retained only as a navigation pointer; every load-bearing step is inside the delivered span.\par\medskip
\ifdefined\hypertarget\hypertarget{sourcepointer-3}{}\fi
\noindent\textit{Source reference.} Original lemma 01ZN: sources/stacks-chow/tags.txt, source lines 2600--2600. Cross-chapter standard prerequisite located in the project tag list staged in the local snapshot; invoked at proof line 10801 of the counted proof. It is retained only as a navigation pointer; every load-bearing step is inside the delivered span.\par\medskip
\ifdefined\hypertarget\hypertarget{sourcepointer-4}{}\fi
\noindent\textit{Source reference.} Original lemma 0ETU: sources/stacks-flat/pinned-original-source.tex, source lines 10303--10303. Same-chapter pointer into the author environment lemma at source line 10303; invoked at proof line 10810 of the counted proof. It is retained only as a navigation pointer; every load-bearing step is inside the delivered span.\par\medskip
\ifdefined\hypertarget\hypertarget{sourcepointer-5}{}\fi
\noindent\textit{Source reference.} Original lemma 081F: sources/stacks-chow/tags.txt, source lines 10215--10215. Cross-chapter standard prerequisite located in the project tag list staged in the local snapshot; invoked at proof line 10817 of the counted proof. It is retained only as a navigation pointer; every load-bearing step is inside the delivered span.\par\medskip
\ifdefined\hypertarget\hypertarget{sourcepointer-6}{}\fi
\noindent\textit{Source reference.} Original lemma 07RR: sources/stacks-chow/tags.txt, source lines 9876--9876. Cross-chapter standard prerequisite located in the project tag list staged in the local snapshot; invoked at proof line 10818 of the counted proof. It is retained only as a navigation pointer; every load-bearing step is inside the delivered span.\par\medskip
\ifdefined\hypertarget\hypertarget{sourcepointer-7}{}\fi
\noindent\textit{Source reference.} Original lemma 01Z4: sources/stacks-chow/tags.txt, source lines 2579--2579. Cross-chapter standard prerequisite located in the project tag list staged in the local snapshot; invoked at proof line 10819 of the counted proof. It is retained only as a navigation pointer; every load-bearing step is inside the delivered span.\par\medskip
\begin{lemma}
\label{lemma-limit-h-topology}
Let $T$ be an affine scheme which is written as a limit
$T = \lim_{i \in I} T_i$ of a directed inverse system of affine schemes.
\begin{enumerate}
\item Let $\mathcal{V} = \{V_j \to T\}_{j = 1, \ldots, m}$ be a
standard h covering of $T$, see Definition
\ref{definition-standard-h}.
Then there exists an index $i$ and a standard h covering
$\mathcal{V}_i = \{V_{i, j} \to T_i\}_{j = 1, \ldots, m}$
whose base change $T \times_{T_i} \mathcal{V}_i$ to $T$
is isomorphic to $\mathcal{V}$.
\item Let $\mathcal{V}_i$, $\mathcal{V}'_i$ be a pair of standard
h coverings of $T_i$. If
$f : T \times_{T_i} \mathcal{V}_i \to T \times_{T_i} \mathcal{V}'_i$ is
a morphism of coverings of $T$, then there exists an index
$i' \geq i$ and a morphism
$f_{i'} : T_{i'} \times_{T_i} \mathcal{V} \to
T_{i'} \times_{T_i} \mathcal{V}'_i$
whose base change to $T$ is $f$.
\item If
$f, g : \mathcal{V} \to \mathcal{V}'_i$
are morphisms of standard h coverings of $T_i$ whose
base changes $f_T, g_T$ to $T$ are equal then there exists an
index $i' \geq i$ such that $f_{T_{i'}} = g_{T_{i'}}$.
\end{enumerate}
In other words, the category of standard h coverings of $T$ is
the colimit over $I$ of the categories of standard h coverings of $T_i$.
\end{lemma}

\begin{proof}
By Limits, Lemma \ref{limits-lemma-descend-finite-presentation}
the category of schemes of finite presentation over $T$ is the
colimit over $I$ of the categories of finite presentation over $T_i$. By
Limits, Lemma \ref{limits-lemma-descend-affine-finite-presentation}
the same is true for category of schemes which are affine and
of finite presentation over $T$.
To finish the proof of the lemma it suffices to show that if
$\{V_{j, i} \to T_i\}_{j = 1, \ldots, m}$ is a finite family of
finitely presented morphisms with $V_{j, i}$ affine, and the
base change family $\{T \times_{T_i} V_{j, i} \to T\}$ is
an h covering, then for some $i' \geq i$ the family
$\{T_{i'} \times_{T_i} V_{j, i} \to T_{i'}\}$ is an h covering.
To see this we use Lemma \ref{lemma-approximate-h-cover} to
choose a finitely presented, proper, surjective
morphism $Y \to T$ and a finite affine open covering
$Y = \bigcup_{k = 1, \ldots, n} Y_k$ such that
$\{Y_k \to T\}_{k = 1, \ldots, n}$ refines
$\{T \times_{T_i} V_{j, i} \to T\}$.
Using the arguments above and
Limits, Lemmas \ref{limits-lemma-eventually-proper},
\ref{limits-lemma-descend-surjective}, and
\ref{limits-lemma-descend-opens}
we can find an $i' \geq i$ and a finitely presented, surjective, proper
morphism $Y_{i'} \to T_{i'}$ and an affine open covering
$Y_{i'} = \bigcup_{k = 1, \ldots, n} Y_{i', k}$
such that moreover $\{Y_{i', k} \to Y_{i'}\}$ refines
$\{T_{i'} \times_{T_i} V_{j, i} \to T_{i'}\}$.
It follows that this last mentioned family is
a h covering and the proof is complete.
\end{proof}

\end{document}
